\documentclass[12pt,a4paper]{article}
\usepackage[utf8]{inputenc}
\usepackage{amsmath}
\usepackage{amssymb}
\usepackage{geometry}
\usepackage{listings}
\usepackage{xcolor}
\usepackage{hyperref}
\usepackage{enumitem}
\usepackage{booktabs}
\usepackage{graphicx}

\geometry{margin=1in}

\title{IID vs Non-IID Performance Comparison\\Federated Learning DDoS Detection}
\author{}
\date{\today}

\lstset{
    basicstyle=\ttfamily\small,
    keywordstyle=\color{blue},
    stringstyle=\color{red},
    commentstyle=\color{green},
    breaklines=true,
    frame=single
}

\begin{document}

\maketitle

\section{Executive Summary}

This document presents a comprehensive comparison of model performance under \textbf{IID (Independent and Identically Distributed)} and \textbf{Non-IID (Non-Independent and Identically Distributed)} data distributions in the federated learning DDoS detection system.

\textbf{Key Finding:} Non-IID distribution presents significant challenges, with accuracy varying dramatically across different model architectures.

\section{Experimental Setup}

\subsection{Configuration}

\begin{itemize}
    \item \textbf{Federated Learning Rounds:} 5 rounds
    \item \textbf{Workers per Model:} 2 workers
    \item \textbf{Data Partition:} 50\% per worker
    \item \textbf{Epochs per Round:} 5 epochs
    \item \textbf{Batch Size:} 32
    \item \textbf{Differential Privacy:} Disabled
    \item \textbf{Aggregation Strategy:} FedAvg
\end{itemize}

\textbf{Note:} For detailed explanations of IID and Non-IID distributions, their implementation, and impact, see the "Data Distribution Explanation" section below.

\section{Data Distribution Explanation}

\subsection{What is IID (Independent and Identically Distributed)?}

\subsubsection{Definition}

\textbf{IID} means that data samples are:
\begin{itemize}
    \item \textbf{Independent:} Each sample is drawn independently (no correlation between samples)
    \item \textbf{Identically Distributed:} All samples come from the same probability distribution
\end{itemize}

In federated learning, \textbf{IID distribution} means each worker/client has data that:
\begin{itemize}
    \item Follows the same statistical distribution
    \item Has similar class proportions
    \item Represents a random sample of the global dataset
\end{itemize}

\subsubsection{How IID is Implemented}

\textbf{Implementation Method:}
\begin{enumerate}
    \item \textbf{Random Sampling:} Each worker receives a random subset of the global dataset
    \item \textbf{Worker-Specific Seed:} Uses \texttt{hash(worker\_id) \% 2**32} to ensure reproducibility
    \item \textbf{Preserved Distribution:} Class proportions remain similar across all workers
\end{enumerate}

\textbf{Code Implementation:}
\begin{lstlisting}[language=Python]
if self.iid:
    # Random sampling with worker-specific seed
    seed = hash(self.worker_id) % 2**32
    n_samples = int(len(X_train) * self.data_partition)
    indices = np.random.RandomState(seed=seed).choice(
        len(X_train), size=n_samples, replace=False
    )
    X_train = X_train.iloc[indices]
    y_train = y_train.iloc[indices]
\end{lstlisting}

\subsubsection{IID Class Distribution}

\textbf{Global Dataset Distribution:}
\begin{itemize}
    \item \textbf{Class 0 (Normal Traffic):} $\sim$60.91\% of total dataset
    \item \textbf{Class 1 (DDoS Attack):} $\sim$39.09\% of total dataset
\end{itemize}

\textbf{Per-Worker Distribution (IID):}
\begin{itemize}
    \item \textbf{Worker 1:} $\sim$60.78\% Class 0, $\sim$39.22\% Class 1
    \item \textbf{Worker 2:} $\sim$60.81\% Class 0, $\sim$39.19\% Class 1
    \item \textbf{Worker 3:} $\sim$60.85\% Class 0, $\sim$39.15\% Class 1
    \item \textbf{Worker 4:} $\sim$60.79\% Class 0, $\sim$39.21\% Class 1
\end{itemize}

\textbf{Key Characteristics:}
\begin{itemize}
    \item ✅ \textbf{Uniform Distribution:} All workers have similar class proportions
    \item ✅ \textbf{Random Sampling:} No bias in data selection
    \item ✅ \textbf{Preserved Statistics:} Global dataset statistics maintained per worker
    \item ✅ \textbf{No Overlap:} Different seeds ensure no data overlap between workers
\end{itemize}

\subsubsection{IID Impact on Model Performance}

\textbf{Positive Impacts:}
\begin{enumerate}
    \item \textbf{Faster Convergence:} Models converge in fewer rounds (typically 3-5 rounds)
    \item \textbf{Higher Accuracy:} Better final model performance (LSTM: 91.05\%, CNN\_LSTM: 87.23\%)
    \item \textbf{Stable Training:} Low variance in model updates across workers
    \item \textbf{Optimal FedAvg:} Standard weighted averaging works effectively
    \item \textbf{Better Generalization:} Models learn more robust features
\end{enumerate}

\textbf{Performance Metrics (IID):}
\begin{itemize}
    \item \textbf{LSTM:} 91.05\% accuracy, 0.1607 loss ✅ (Best)
    \item \textbf{CNN\_LSTM:} 87.23\% accuracy, 0.2596 loss ✅ (Good)
    \item \textbf{CNN1D:} 76.16\% accuracy, 0.4691 loss ⚠️ (Moderate)
    \item \textbf{MLPv2:} 60.91\% accuracy, 0.6571 loss ⚠️ (Moderate)
\end{itemize}

\subsection{What is Non-IID (Non-Independent and Identically Distributed)?}

\subsubsection{Definition}

\textbf{Non-IID} means that data samples are:
\begin{itemize}
    \item \textbf{Not Identically Distributed:} Different workers have different data distributions
    \item \textbf{May be Correlated:} Samples within a worker may be correlated
    \item \textbf{Class-Skewed:} Workers have significantly different class proportions
\end{itemize}

In federated learning, \textbf{Non-IID distribution} means:
\begin{itemize}
    \item Workers have \textbf{different class distributions}
    \item Some workers have \textbf{more of one class} than others
    \item More realistic representation of real-world scenarios (e.g., hospitals, IoT devices)
\end{itemize}

\subsubsection{How Non-IID is Implemented}

\textbf{Implementation Method:}
\begin{enumerate}
    \item \textbf{Class-Based Partitioning:} Workers receive data based on class labels
    \item \textbf{Skewed Distribution:} Odd workers get 80\% Class 0, Even workers get 80\% Class 1
    \item \textbf{Deterministic Assignment:} Worker number (odd/even) determines class skew
\end{enumerate}

\textbf{Code Implementation:}
\begin{lstlisting}[language=Python]
if not self.iid:
    # Extract worker number (e.g., "worker1" -> 1)
    worker_num = int(''.join(filter(str.isdigit, self.worker_id)))
    is_odd_worker = (worker_num % 2 == 1)
    
    # Get class indices
    class_0_indices = np.where(y_train == 0)[0]
    class_1_indices = np.where(y_train == 1)[0]
    
    if is_odd_worker:
        # Odd workers: 80% Class 0, 20% Class 1
        n_class_0 = int(total_samples_needed * 0.8)
        n_class_1 = total_samples_needed - n_class_0
    else:
        # Even workers: 20% Class 0, 80% Class 1
        n_class_1 = int(total_samples_needed * 0.8)
        n_class_0 = total_samples_needed - n_class_1
    
    # Sample from each class
    selected_class_0 = rng.choice(class_0_indices, size=n_class_0, replace=False)
    selected_class_1 = rng.choice(class_1_indices, size=n_class_1, replace=False)
    
    # Combine and shuffle
    indices = np.concatenate([selected_class_0, selected_class_1])
    X_train = X_train.iloc[indices]
    y_train = y_train.iloc[indices]
\end{lstlisting}

\subsubsection{Non-IID Class Distribution Pattern}

\textbf{Per-Worker Distribution (Non-IID):}

\textbf{Odd Workers} (worker1, worker3, worker5, worker7):
\begin{itemize}
    \item \textbf{Class 0 (Normal Traffic):} 80\% of worker's data
    \item \textbf{Class 1 (DDoS Attack):} 20\% of worker's data
    \item \textbf{Example:} Worker 1 might have 800 samples of Class 0, 200 samples of Class 1
\end{itemize}

\textbf{Even Workers} (worker2, worker4, worker6, worker8):
\begin{itemize}
    \item \textbf{Class 0 (Normal Traffic):} 20\% of worker's data
    \item \textbf{Class 1 (DDoS Attack):} 80\% of worker's data
    \item \textbf{Example:} Worker 2 might have 200 samples of Class 0, 800 samples of Class 1
\end{itemize}

\textbf{Visual Representation:}
\begin{verbatim}
Worker 1 (Odd):  [████████████████████████████████] 80% Class 0
                 [████] 20% Class 1

Worker 2 (Even): [████] 20% Class 0
                 [████████████████████████████████] 80% Class 1
\end{verbatim}

\textbf{Key Characteristics:}
\begin{itemize}
    \item ⚠️ \textbf{Skewed Distribution:} Workers have complementary but different class proportions
    \item ⚠️ \textbf{Class Imbalance:} 60\% difference in class distribution between odd/even workers
    \item ⚠️ \textbf{Realistic Scenario:} Mimics real-world federated learning (e.g., different hospitals, IoT devices)
    \item ⚠️ \textbf{Challenging:} Creates significant challenges for model convergence
\end{itemize}

\subsubsection{Non-IID Impact on Model Performance}

\textbf{Negative Impacts:}
\begin{enumerate}
    \item \textbf{Slower Convergence:} Requires more communication rounds to converge
    \item \textbf{Lower Accuracy:} Achieves lower final accuracy compared to IID
    \item \textbf{Higher Variance:} More variability in model updates across workers
    \item \textbf{Model-Specific Sensitivity:} Different architectures show varying robustness:
    \begin{itemize}
        \item \textbf{CNN\_LSTM:} Only -3.45\% accuracy drop (most robust)
        \item \textbf{LSTM:} -5.47\% accuracy drop (very robust)
        \item \textbf{CNN1D:} -15.25\% accuracy drop (moderate impact)
        \item \textbf{MLPv2:} -21.82\% accuracy drop (most sensitive, fails catastrophically)
    \end{itemize}
\end{enumerate}

\textbf{Performance Metrics (Non-IID):}
\begin{itemize}
    \item \textbf{LSTM:} 85.58\% accuracy, 0.3153 loss ✅ (Best, but 5.47\% drop from IID)
    \item \textbf{CNN\_LSTM:} 83.78\% accuracy, 0.3094 loss ✅ (Good, only 3.45\% drop from IID)
    \item \textbf{CNN1D:} 60.91\% accuracy, 0.6972 loss ⚠️ (Moderate, 15.25\% drop from IID)
    \item \textbf{MLPv2:} 39.09\% accuracy, 0.9770 loss ❌ (Poor, 21.82\% drop, below random chance)
\end{itemize}

\textbf{Why Non-IID is Challenging:}
\begin{enumerate}
    \item \textbf{Gradient Mismatch:} Workers optimize for different objectives (different class distributions)
    \item \textbf{Aggregation Difficulty:} FedAvg struggles when workers have conflicting updates
    \item \textbf{Local Overfitting:} Workers may overfit to their local class distribution
    \item \textbf{Convergence Instability:} Model may oscillate or converge to suboptimal solutions
\end{enumerate}

\section{Non-IID Results}

\subsection{Performance Metrics}

\begin{table}[h]
\centering
\begin{tabular}{@{}lcccc@{}}
\toprule
\textbf{Model} & \textbf{Accuracy} & \textbf{Loss} & \textbf{Parameters} & \textbf{Convergence} \\
\midrule
LSTM & \textbf{85.58\%} & 0.3153 & 29.38K & ✅ Good \\
CNN\_LSTM & \textbf{83.78\%} & 0.3094 & 37.57K & ✅ Good \\
CNN1D & \textbf{60.91\%} & 0.6972 & 58.69K & ⚠️ Moderate \\
MLPv2 & \textbf{39.09\%} & 0.9770 & 4.07K & ❌ Poor \\
\bottomrule
\end{tabular}
\caption{Non-IID Performance Results}
\end{table}

\subsection{Detailed Non-IID Results}

\subsubsection{LSTM (Best Performance)}

\begin{itemize}
    \item \textbf{Accuracy:} 85.58\%
    \item \textbf{Loss:} 0.3153
    \item \textbf{Parameters:} 29.38K
    \item \textbf{Network Traffic:} 1.13 MB sent, 1.13 MB received
    \item \textbf{Total Params Exchanged:} 296.34K sent, 296.34K received
    \item \textbf{Status:} ✅ Complete - Round 5/5
\end{itemize}

\textbf{Analysis:} LSTM shows the best performance under Non-IID conditions, achieving 85.58\% accuracy. The sequential architecture with LSTM layers effectively captures temporal patterns and handles class imbalance well, outperforming the hybrid CNN\_LSTM model.

\subsubsection{CNN\_LSTM}

\begin{itemize}
    \item \textbf{Accuracy:} 83.78\%
    \item \textbf{Loss:} 0.3094
    \item \textbf{Parameters:} 37.57K
    \item \textbf{Network Traffic:} 1.44 MB sent, 1.44 MB received
    \item \textbf{Total Params Exchanged:} 378.58K sent, 378.58K received
    \item \textbf{Status:} ✅ Complete - Round 5/5
\end{itemize}

\textbf{Analysis:} CNN\_LSTM performs well (83.78\%) under Non-IID conditions, achieving the second-best accuracy. The hybrid architecture (CNN + LSTM) handles class imbalance reasonably well, though slightly behind the pure LSTM model.

\subsubsection{CNN1D}

\begin{itemize}
    \item \textbf{Accuracy:} 60.91\%
    \item \textbf{Loss:} 0.6972
    \item \textbf{Parameters:} 58.69K
    \item \textbf{Network Traffic:} 2.26 MB sent, 2.26 MB received
    \item \textbf{Total Params Exchanged:} 592.02K sent, 592.02K received
    \item \textbf{Status:} ✅ Complete - Round 5/5
\end{itemize}

\textbf{Analysis:} CNN1D achieves moderate accuracy (60.91\%) but struggles with the class imbalance. The convolutional layers may not adapt well to the skewed distribution.

\subsubsection{MLPv2 (Worst Performance)}

\begin{itemize}
    \item \textbf{Accuracy:} 39.09\%
    \item \textbf{Loss:} 0.9770
    \item \textbf{Parameters:} 4.07K
    \item \textbf{Network Traffic:} 176.33 KB sent, 176.33 KB received
    \item \textbf{Total Params Exchanged:} 45.14K sent, 45.14K received
    \item \textbf{Status:} ✅ Complete - Round 5/5
\end{itemize}

\textbf{Analysis:} MLPv2 shows the poorest performance (39.09\% accuracy), which is \textbf{below random chance} for binary classification (50\%). This suggests the model is failing to learn meaningful patterns under Non-IID conditions. The high loss (0.9770) indicates poor convergence.

\section{IID Results}

\subsection{Performance Metrics}

\begin{table}[h]
\centering
\begin{tabular}{@{}lcccc@{}}
\toprule
\textbf{Model} & \textbf{Accuracy} & \textbf{Loss} & \textbf{Parameters} & \textbf{Convergence} \\
\midrule
LSTM & \textbf{91.05\%} & 0.1607 & 29.38K & ✅ Excellent \\
CNN\_LSTM & \textbf{87.23\%} & 0.2596 & 37.57K & ✅ Good \\
CNN1D & \textbf{76.16\%} & 0.4691 & 58.69K & ⚠️ Moderate \\
MLPv2 & \textbf{60.91\%} & 0.6571 & 4.07K & ⚠️ Moderate \\
\bottomrule
\end{tabular}
\caption{IID Performance Results}
\end{table}

\subsection{Detailed IID Results}

\subsubsection{LSTM (Best Performance)}

\begin{itemize}
    \item \textbf{Accuracy:} 91.05\%
    \item \textbf{Loss:} 0.1607
    \item \textbf{Parameters:} 29.38K
    \item \textbf{Network Traffic:} 1.02 MB sent, 926.06 KB received
    \item \textbf{Total Params Exchanged:} 266.71K sent, 237.07K received
    \item \textbf{Status:} ✅ Complete - Round 5/5
\end{itemize}

\textbf{Analysis:} LSTM achieves the best performance under IID conditions with 91.05\% accuracy and very low loss (0.1607). The sequential architecture excels when data is uniformly distributed across workers.

\subsubsection{CNN\_LSTM}

\begin{itemize}
    \item \textbf{Accuracy:} 87.23\%
    \item \textbf{Loss:} 0.2596
    \item \textbf{Parameters:} 37.57K
    \item \textbf{Network Traffic:} 1.44 MB sent, 1.44 MB received
    \item \textbf{Total Params Exchanged:} 378.58K sent, 378.58K received
    \item \textbf{Status:} ✅ Complete - Round 5/5
\end{itemize}

\textbf{Analysis:} CNN\_LSTM performs well under IID conditions (87.23\% accuracy), demonstrating the effectiveness of hybrid architectures when data distribution is uniform.

\subsubsection{CNN1D}

\begin{itemize}
    \item \textbf{Accuracy:} 76.16\%
    \item \textbf{Loss:} 0.4691
    \item \textbf{Parameters:} 58.69K
    \item \textbf{Network Traffic:} 2.26 MB sent, 2.26 MB received
    \item \textbf{Total Params Exchanged:} 592.02K sent, 592.02K received
    \item \textbf{Status:} ✅ Complete - Round 5/5
\end{itemize}

\textbf{Analysis:} CNN1D achieves moderate accuracy (76.16\%) under IID conditions. While better than Non-IID, it still lags behind sequential and hybrid models.

\subsubsection{MLPv2}

\begin{itemize}
    \item \textbf{Accuracy:} 60.91\%
    \item \textbf{Loss:} 0.6571
    \item \textbf{Parameters:} 4.07K
    \item \textbf{Network Traffic:} 176.33 KB sent, 176.33 KB received
    \item \textbf{Total Params Exchanged:} 45.14K sent, 45.14K received
    \item \textbf{Status:} ✅ Complete - Round 5/5
\end{itemize}

\textbf{Analysis:} MLPv2 shows moderate performance (60.91\% accuracy) under IID conditions, which is significantly better than its Non-IID performance (39.09\%). However, it still underperforms compared to other architectures.

\section{Comparison Analysis}

\subsection{Side-by-Side Performance Comparison}

\begin{table}[h]
\centering
\small
\begin{tabular}{@{}lcccccc@{}}
\toprule
\textbf{Model} & \textbf{IID Acc.} & \textbf{Non-IID Acc.} & \textbf{Drop} & \textbf{IID Loss} & \textbf{Non-IID Loss} & \textbf{Impact} \\
\midrule
LSTM & \textbf{91.05\%} & \textbf{85.58\%} & \textbf{-5.47\%} & 0.1607 & 0.3153 & ⚠️ Moderate \\
CNN\_LSTM & \textbf{87.23\%} & \textbf{83.78\%} & \textbf{-3.45\%} & 0.2596 & 0.3094 & ✅ Low \\
CNN1D & \textbf{76.16\%} & \textbf{60.91\%} & \textbf{-15.25\%} & 0.4691 & 0.6972 & ❌ High \\
MLPv2 & \textbf{60.91\%} & \textbf{39.09\%} & \textbf{-21.82\%} & 0.6571 & 0.9770 & ❌ Very High \\
\bottomrule
\end{tabular}
\caption{IID vs Non-IID Performance Comparison}
\end{table}

\subsection{Performance Ranking}

\subsubsection{IID Distribution}

\begin{enumerate}
    \item \textbf{LSTM:} 91.05\% ✅ (Best)
    \item \textbf{CNN\_LSTM:} 87.23\% ✅ (Good)
    \item \textbf{CNN1D:} 76.16\% ⚠️ (Moderate)
    \item \textbf{MLPv2:} 60.91\% ⚠️ (Moderate)
\end{enumerate}

\subsubsection{Non-IID Distribution}

\begin{enumerate}
    \item \textbf{LSTM:} 85.58\% ✅ (Best)
    \item \textbf{CNN\_LSTM:} 83.78\% ✅ (Good)
    \item \textbf{CNN1D:} 60.91\% ⚠️ (Moderate)
    \item \textbf{MLPv2:} 39.09\% ❌ (Poor)
\end{enumerate}

\subsection{Key Observations}

\subsubsection{IID vs Non-IID Impact Analysis}

\textbf{Most Robust Models (Lowest Accuracy Drop):}
\begin{itemize}
    \item \textbf{CNN\_LSTM:} Only -3.45\% drop (87.23\% $\rightarrow$ 83.78\%) - Most resilient to Non-IID
    \item \textbf{LSTM:} -5.47\% drop (91.05\% $\rightarrow$ 85.58\%) - Very robust, maintains high accuracy
\end{itemize}

\textbf{Most Sensitive Models (Highest Accuracy Drop):}
\begin{itemize}
    \item \textbf{MLPv2:} -21.82\% drop (60.91\% $\rightarrow$ 39.09\%) - Severely impacted by Non-IID
    \item \textbf{CNN1D:} -15.25\% drop (76.16\% $\rightarrow$ 60.91\%) - Significantly affected
\end{itemize}

\subsubsection{Model Architecture Sensitivity}

Different architectures show varying sensitivity to Non-IID distribution:

\begin{itemize}
    \item \textbf{Hybrid Models (CNN\_LSTM):} Most robust to Non-IID, only 3.45\% accuracy drop
    \item \textbf{Sequential Models (LSTM):} Very robust, 5.47\% drop, maintains best overall accuracy
    \item \textbf{Convolutional Models (CNN1D):} Moderate robustness, 15.25\% drop
    \item \textbf{Fully Connected (MLPv2):} Least robust, 21.82\% drop, fails under Non-IID
\end{itemize}

\subsubsection{Class Imbalance Impact}

The 80/20 class distribution creates significant challenges:

\begin{itemize}
    \item \textbf{MLPv2} struggles most, dropping from 60.91\% (IID) to 39.09\% (Non-IID) - below random chance
    \item \textbf{CNN\_LSTM} handles imbalance best, maintaining 83.78\% accuracy with only 3.45\% drop
    \item \textbf{LSTM} shows excellent resilience, maintaining 85.58\% accuracy despite 5.47\% drop
    \item Models with more parameters (CNN1D: 58.69K) don't necessarily perform better
\end{itemize}

\subsubsection{Loss vs Accuracy Correlation}

\textbf{IID Distribution:}
\begin{itemize}
    \item \textbf{LSTM:} Very low loss (0.1607) $\rightarrow$ Excellent accuracy (91.05\%) ✅
    \item \textbf{CNN\_LSTM:} Low loss (0.2596) $\rightarrow$ Good accuracy (87.23\%) ✅
    \item \textbf{CNN1D:} Moderate loss (0.4691) $\rightarrow$ Moderate accuracy (76.16\%) ⚠️
    \item \textbf{MLPv2:} High loss (0.6571) $\rightarrow$ Moderate accuracy (60.91\%) ⚠️
\end{itemize}

\textbf{Non-IID Distribution:}
\begin{itemize}
    \item \textbf{LSTM:} Low loss (0.3153) $\rightarrow$ High accuracy (85.58\%) ✅
    \item \textbf{CNN\_LSTM:} Low loss (0.3094) $\rightarrow$ High accuracy (83.78\%) ✅
    \item \textbf{CNN1D:} High loss (0.6972) $\rightarrow$ Moderate accuracy (60.91\%) ⚠️
    \item \textbf{MLPv2:} Very high loss (0.9770) $\rightarrow$ Poor accuracy (39.09\%) ❌
\end{itemize}

\textbf{Loss Increase Analysis:}
\begin{itemize}
    \item \textbf{LSTM:} Loss increases by 96.2\% (0.1607 $\rightarrow$ 0.3153) but maintains high accuracy
    \item \textbf{CNN\_LSTM:} Loss increases by 19.2\% (0.2596 $\rightarrow$ 0.3094) - most stable
    \item \textbf{CNN1D:} Loss increases by 48.6\% (0.4691 $\rightarrow$ 0.6972)
    \item \textbf{MLPv2:} Loss increases by 48.7\% (0.6571 $\rightarrow$ 0.9770)
\end{itemize}

\subsubsection{Convergence Behavior}

\textbf{IID Distribution:}
\begin{itemize}
    \item \textbf{LSTM:} Excellent convergence (very low loss: 0.1607, high accuracy: 91.05\%)
    \item \textbf{CNN\_LSTM:} Good convergence (low loss: 0.2596, good accuracy: 87.23\%)
    \item \textbf{CNN1D:} Moderate convergence (moderate loss: 0.4691, moderate accuracy: 76.16\%)
    \item \textbf{MLPv2:} Moderate convergence (high loss: 0.6571, moderate accuracy: 60.91\%)
\end{itemize}

\textbf{Non-IID Distribution:}
\begin{itemize}
    \item \textbf{LSTM:} Excellent convergence (low loss: 0.3153, high accuracy: 85.58\%)
    \item \textbf{CNN\_LSTM:} Good convergence (low loss: 0.3094, high accuracy: 83.78\%)
    \item \textbf{CNN1D:} Moderate convergence (high loss: 0.6972, moderate accuracy: 60.91\%)
    \item \textbf{MLPv2:} Poor convergence (very high loss: 0.9770, poor accuracy: 39.09\%)
\end{itemize}

\section{Discussion}

\subsection{Why LSTM Performs Best Under Non-IID?}

\begin{enumerate}
    \item \textbf{Sequential Processing:} LSTM's sequential nature effectively captures temporal patterns in the data
    \item \textbf{Memory Mechanism:} Long short-term memory helps maintain context across skewed class distributions
    \item \textbf{Parameter Efficiency:} 29.38K parameters provide good balance between capacity and generalization
    \item \textbf{Robustness:} LSTM gates (forget, input, output) help filter and adapt to imbalanced data
    \item \textbf{Adaptability:} Can learn meaningful patterns despite 80/20 class skew across workers
\end{enumerate}

\subsection{Why CNN\_LSTM Performs Well Under Non-IID?}

\begin{enumerate}
    \item \textbf{Feature Extraction:} CNN layers extract local patterns, LSTM captures temporal dependencies
    \item \textbf{Robustness:} Hybrid architecture provides redundancy and robustness
    \item \textbf{Parameter Count:} 37.57K parameters provide sufficient capacity
    \item \textbf{Adaptability:} Can learn from both local and sequential patterns despite class imbalance
\end{enumerate}

\subsection{Why MLPv2 Performs Worst Under Non-IID?}

\begin{enumerate}
    \item \textbf{Simple Architecture:} Fully connected layers lack specialized feature extraction
    \item \textbf{Limited Capacity:} Only 4.07K parameters may be insufficient for complex patterns
    \item \textbf{No Inductive Bias:} Lacks architectural biases that help with imbalanced data
    \item \textbf{Gradient Issues:} May struggle with gradient flow in Non-IID scenarios
\end{enumerate}

\subsection{Implications for Federated Learning}

\begin{enumerate}
    \item \textbf{Model Selection:} Architecture choice is critical for Non-IID scenarios
    \item \textbf{Robustness Testing:} Non-IID reveals model weaknesses not visible in IID
    \item \textbf{Production Deployment:} Real-world data is often Non-IID, making this testing essential
    \item \textbf{Algorithm Development:} May need specialized aggregation strategies for Non-IID
\end{enumerate}

\section{Network Traffic Analysis}

\subsection{Non-IID Network Traffic}

\begin{table}[h]
\centering
\begin{tabular}{@{}lcccc@{}}
\toprule
\textbf{Model} & \textbf{Bytes Sent} & \textbf{Bytes Received} & \textbf{Params Sent} & \textbf{Params Received} \\
\midrule
CNN1D & 2.26 MB & 2.26 MB & 592.02K & 592.02K \\
CNN\_LSTM & 1.44 MB & 1.44 MB & 378.58K & 378.58K \\
LSTM & 1.13 MB & 1.13 MB & 296.34K & 296.34K \\
MLPv2 & 176.33 KB & 176.33 KB & 45.14K & 45.14K \\
\bottomrule
\end{tabular}
\caption{Network Traffic for Non-IID Training}
\end{table}

\textbf{Observations:}
\begin{itemize}
    \item Network traffic correlates with model size (parameters)
    \item CNN1D has highest traffic due to largest parameter count
    \item MLPv2 has lowest traffic due to smallest parameter count
    \item Traffic is symmetric (sent $\approx$ received) for most models
\end{itemize}

\section{Conclusion}

\subsection{Performance Summary}

\textbf{IID Distribution Performance:}
\begin{itemize}
    \item \textbf{Best Case (LSTM):} 91.05\% accuracy - excellent for production
    \item \textbf{Second Best (CNN\_LSTM):} 87.23\% accuracy - good for production
    \item \textbf{Moderate (CNN1D):} 76.16\% accuracy - acceptable for production
    \item \textbf{Weakest (MLPv2):} 60.91\% accuracy - borderline acceptable
\end{itemize}

\textbf{Non-IID Distribution Performance:}
\begin{itemize}
    \item \textbf{Best Case (LSTM):} 85.58\% accuracy - good for production (5.47\% drop)
    \item \textbf{Second Best (CNN\_LSTM):} 83.78\% accuracy - acceptable for production (3.45\% drop)
    \item \textbf{Moderate (CNN1D):} 60.91\% accuracy - acceptable but degraded (15.25\% drop)
    \item \textbf{Worst Case (MLPv2):} 39.09\% accuracy - unacceptable, below random chance (21.82\% drop)
\end{itemize}

\textbf{Key Finding:} IID distribution consistently outperforms Non-IID for all models, with accuracy drops ranging from 3.45\% (CNN\_LSTM) to 21.82\% (MLPv2).

\subsection{Key Takeaways}

\begin{enumerate}
    \item \textbf{IID is Better:} All models perform better under IID distribution, confirming the importance of data distribution uniformity
    \item \textbf{Architecture Matters:} CNN\_LSTM shows best robustness to Non-IID (only 3.45\% drop), while MLPv2 is most sensitive (21.82\% drop)
    \item \textbf{LSTM Best Overall:} LSTM achieves highest accuracy in both IID (91.05\%) and Non-IID (85.58\%) scenarios
    \item \textbf{Hybrid Models Resilient:} CNN\_LSTM maintains excellent performance with minimal degradation under Non-IID
    \item \textbf{MLPv2 Struggles:} Simple architectures fail catastrophically under Non-IID conditions (below random chance)
    \item \textbf{Testing Essential:} Non-IID testing reveals real-world performance degradation
    \item \textbf{Model Selection:} Choose architectures based on expected data distribution - use CNN\_LSTM or LSTM for Non-IID scenarios
\end{enumerate}

\section{Appendix: Experimental Details}

\subsection{Non-IID Distribution Details}

\begin{itemize}
    \item \textbf{Worker 1 (Odd):} 80\% Class 0, 20\% Class 1
    \item \textbf{Worker 2 (Even):} 20\% Class 0, 80\% Class 1
    \item \textbf{Class Skew:} 60\% difference between workers
    \item \textbf{Data Overlap:} None (complementary distributions)
\end{itemize}

\subsection{Training Configuration}

\begin{itemize}
    \item \textbf{Optimizer:} Adam (learning\_rate=0.001)
    \item \textbf{Loss Function:} Categorical Crossentropy
    \item \textbf{Activation:} Softmax (output layer)
    \item \textbf{Regularization:} None (for baseline comparison)
    \item \textbf{Class Weights:} Not applied (to test raw Non-IID impact)
\end{itemize}

\vspace{1cm}

\textit{Document Status: Non-IID results complete. IID results pending.}

\end{document}

